\documentclass[aspectratio=169,11pt]{beamer}

% ── Session 3: "Drawing Inside the Labyrinth" (fractal encoding) ──────────
% Shared brand theme + \gwtitle/\punch/\memeslot live in gwtheme.tex.

\usepackage[utf8]{inputenc}
\usepackage[T1]{fontenc}
\usepackage{lmodern}
\usepackage{../style/greatwall}
\usepackage{booktabs}
\usepackage{array}
\usepackage{colortbl}

\graphicspath{{./}{../assets/}}
\def\gwsessionno{3}
% gwtheme.tex — shared Beamer theme for the Great Wall pilot-course decks.
% \input this from each slides/sessionN.tex *after* loading greatwall.sty and
% *after* defining \gwsessionno (the session number, e.g. \def\gwsessionno{2}).
% Provides the brand look plus three helpers: \gwtitle, \punch, \memeslot.
%
% Depends on: beamer, greatwall.sty (palette + boxes + marks), tikz.

\providecommand{\gwsessionno}{?}

% ── Minimal custom theme (no external font deps; texlive-full only) ────────
\usefonttheme{professionalfonts}
\setbeamertemplate{navigation symbols}{}
\setbeamercolor{normal text}{fg=darktext, bg=white}
\setbeamercolor{structure}{fg=accent}
\setbeamercolor{frametitle}{fg=white, bg=accent}
\setbeamercolor{title}{fg=white}
\setbeamercolor{block title}{fg=white, bg=accentdark}
\setbeamercolor{block body}{fg=darktext, bg=accentlight}
\setbeamercolor{itemize item}{fg=accent}
\setbeamercolor{itemize subitem}{fg=btccol}
\setbeamerfont{frametitle}{series=\bfseries}
\setbeamerfont{title}{series=\bfseries, size=\Huge}
\setbeamertemplate{itemize item}{\textbf{--}}

% Frametitle as a solid green bar
\setbeamertemplate{frametitle}{%
  \nointerlineskip
  \begin{beamercolorbox}[wd=\paperwidth,ht=2.6ex,dp=1.2ex,leftskip=0.4cm,rightskip=0.4cm]{frametitle}%
    \usebeamerfont{frametitle}\insertframetitle%
  \end{beamercolorbox}%
}

% Footline: course | session | frame number
\setbeamercolor{gwfootline}{fg=graytext, bg=white}
\setbeamertemplate{footline}{%
  \leavevmode%
  \hbox{%
  \begin{beamercolorbox}[wd=.5\paperwidth,ht=2.4ex,dp=1ex,leftskip=0.4cm]{gwfootline}%
    \tiny Great Wall --- Pilot Course%
  \end{beamercolorbox}%
  \begin{beamercolorbox}[wd=.5\paperwidth,ht=2.4ex,dp=1ex,rightskip=0.4cm]{gwfootline}%
    \hfill\tiny Session \gwsessionno{} \;\textbullet\; \insertframenumber%
  \end{beamercolorbox}}%
}

% ── Helpers ───────────────────────────────────────────────────────────────
% Big punchy centered statement frame (full-bleed green).
\newcommand{\punch}[1]{%
  \begin{frame}[plain]
  \begin{tikzpicture}[remember picture,overlay]
    \fill[accent] (current page.south west) rectangle (current page.north east);
  \end{tikzpicture}
  \begin{center}\vfill
  {\color{white}\Huge\bfseries #1}\par
  \vfill\end{center}
  \end{frame}}

% Meme slot — shows the image if present, else a labelled placeholder box.
% #1 = filename in ../assets/, #2 = caption/placeholder text.
\newcommand{\memeslot}[2]{%
  \IfFileExists{../assets/#1}{%
    \includegraphics[width=\linewidth,height=0.62\textheight,keepaspectratio]{#1}%
  }{%
    \begin{tcolorbox}[colback=accentlight,colframe=graytext,sharp corners,
      halign=center, valign=center, height=0.55\textheight]
      \color{graytext}\itshape [meme slot: \texttt{#1}]\\[4pt]\small #2
    \end{tcolorbox}%
  }%
}

% Branded title frame. #1 = title, #2 = subtitle.
\newcommand{\gwtitle}[2]{%
  \begin{frame}[plain]
    \begin{tikzpicture}[remember picture,overlay]
      \fill[accent] (current page.south west) rectangle (current page.north east);
      \foreach \i in {1,...,8}{
        \fill[white,opacity=0.08] ([xshift=-\i cm,yshift=\i cm]current page.south east)
          rectangle ++(0.7cm,0.7cm);
      }
    \end{tikzpicture}
    \vfill
    \begin{center}
    {\color{white}\Huge\bfseries #1\par}
    \vspace{6pt}
    {\color{accentlight}\large #2\par}
    \vspace{18pt}
    {\color{accentlight}\small Great Wall \;\textbullet\; Session \gwsessionno{} of 5 \;\textbullet\; hands-on, on Discord\par}
    \end{center}
    \vfill
  \end{frame}}


\begin{document}

\gwtitle{Drawing Inside the Labyrinth}{How a secret becomes an exact point on a fractal --- losslessly}

% ── Recap + hook ──────────────────────────────────────────────────────────
\begin{frame}{The obvious objection}
  We said your secret is a \emph{point on a fractal} you can recognize but not
  describe. Fair pushback:
  \begin{block}{}
    ``If you can't even \emph{say} where it is, how is the mapping \textbf{exact},
    \textbf{lossless}, and \textbf{reversible}? And why a \emph{fractal}?''
  \end{block}
  \vspace{4pt}
  Today we answer all three --- and you'll watch your own words round-trip.
  \begin{center}\footnotesize\color{graytext}
  Check-in: one word in \texttt{\#feedback} --- how did the viewer feel to play with?
  \end{center}
\end{frame}

% ── Entropy grounding ─────────────────────────────────────────────────────
\begin{frame}{First: what is a bit of entropy?}
  \begin{columns}[T]
    \begin{column}{0.62\textwidth}
      \small
      \textbf{A bit of entropy?} One \emph{unpredictable} coin-flip. $n$ bits
      $=2^n$ equally-likely secrets to sift. (Unpredictable is the point: a bit
      that's always 1 stores data but \emph{zero} entropy.)
      \vspace{5pt}

      \textbf{Why 128?} $2^{128}\!\approx\!3.4\times10^{38}$ --- uncrackable by
      anyone, ever. (64 = grindable; 256 = overkill.) Also BIP39's 12-word
      default.
      \vspace{5pt}

      \textbf{Why 32 per point?} A point $=$ 32 left/right bisection choices. Pick
      the leaf at random from $2^{32}$ $\Rightarrow$ 32 bits. So
      \textbf{4 points $=$ 128 bits} ($=$ 12 words $=$ 4 stages).
    \end{column}
    \begin{column}{0.36\textwidth}
      \centering
      \begin{tikzpicture}[font=\scriptsize,
        b/.style={draw=btccol, fill=accentlight, rounded corners=1.5pt,
                  minimum width=2.6cm, minimum height=0.55cm}]
        \node[b] (p1) at (0,0)     {point 1 \;=\; 32};
        \node[b] (p2) at (0,-0.72) {point 2 \;=\; 32};
        \node[b] (p3) at (0,-1.44) {point 3 \;=\; 32};
        \node[b] (p4) at (0,-2.16) {point 4 \;=\; 32};
        \draw[accentdark, thick] (-1.45,-2.6) -- (1.45,-2.6);
        \node[accentdark] at (0,-2.95) {\textbf{= 128 bits}};
      \end{tikzpicture}
    \end{column}
  \end{columns}
  \vspace{2pt}
  \begin{center}\footnotesize\color{graytext}
  The \emph{haystack} (which fractal) isn't part of this count --- that's the
  coercion \emph{cost} (Session 4), not seed entropy.
  \end{center}
\end{frame}

% ── BIP39 words + checksum ────────────────────────────────────────────────
\begin{frame}{Two more: BIP39 words and checksums}
  \begin{columns}[T]
    \begin{column}{0.6\textwidth}
      \small
      \textbf{How does a word encode 11 bits?} The BIP39 wordlist is exactly
      \textbf{2048 words} $=2^{11}$. A word is just an index 0--2047 --- and any
      number in that range is 11 bits. So 12 words $\Rightarrow 12\times11 =
      \textbf{132}$ bits.
      \vspace{6pt}

      \textbf{What are checksum bits?} Of the 132, only 128 are entropy; the last
      \textbf{4 are a checksum} --- the first 4 bits of
      $\mathrm{SHA\text{-}256}$(entropy). \emph{Derived}, not random (zero
      entropy); they let a wallet \textbf{catch typos} --- random 12 words pass
      only $1/16$ of the time.
    \end{column}
    \begin{column}{0.38\textwidth}
      \centering
      \vspace{10pt}
      \resizebox{\linewidth}{!}{%
      \begin{tikzpicture}[font=\scriptsize]
        \node[anchor=south] at (2.5,0.32) {12 words $\times$ 11 = \textbf{132} bits};
        \fill[accent!80] (0,0) rectangle (4.3,0.6);
        \fill[attackcol] (4.3,0) rectangle (5.0,0.6);
        \draw[darktext] (0,0) rectangle (5.0,0.6);
        \node[white] at (2.15,0.3) {\textbf{128 entropy}};
        \node[white] at (4.65,0.3) {\textbf{4}};
        \draw[-latex, accentdark, thick] (2.15,-0.15) -- (2.15,-0.6);
        \node[accentdark, anchor=north, text width=3cm, align=center] at (2.15,-0.6)
          {Great Wall stores\\these (4 points)};
        \draw[-latex, attackcol, thick] (4.65,-0.15) -- (4.65,-0.6);
        \node[attackcol, anchor=north] at (4.9,-0.6) {checksum};
      \end{tikzpicture}}
    \end{column}
  \end{columns}
  \vspace{10pt}
  \begin{center}\footnotesize\color{graytext}
  Great Wall encodes only the \textbf{128 entropy bits} (4 points). The 4 checksum
  bits aren't stored --- they're \emph{recomputed} on decode.
  \end{center}
\end{frame}

% ── Big picture ───────────────────────────────────────────────────────────
\begin{frame}{The trick: a shared codebook nobody stores}
  Both encoder and decoder can \textbf{regenerate the same fractal} from the same
  parameters. Its intricate structure is a codebook neither side saves --- they
  each rebuild it on demand.
  \vspace{6pt}
  \begin{keyidea}[Haystack $\gg$ needle]
    Specifying \emph{which} fractal (the haystack) is made \textbf{far more
    expensive} than specifying \emph{where} in it your point sits (the needle,
    $\approx$ 32 bits). That asymmetry is the security engine --- and it
    compounds per stage (Session 4).
  \end{keyidea}
\end{frame}

% ── Bisection encoding ────────────────────────────────────────────────────
\begin{frame}{Bisection encoding, one bit at a time}
  \begin{columns}[T]
    \begin{column}{0.46\textwidth}
      \footnotesize
      \begin{enumerate}
        \item Find \textbf{islands} (structure) in the current rectangle.
        \item The \emph{fractal} picks a split line (weighted median).
        \item Read one bit as a \textbf{direction}: \\
          vertical $0{=}$left $1{=}$right; horizontal $0{=}$up $1{=}$down.
        \item Contract if needed; recurse on the chosen child.
      \end{enumerate}
      \vspace{2pt}
      After all bits: the \textbf{center of the final rectangle} is your point.
    \end{column}
    \begin{column}{0.52\textwidth}
      \centering
      \begin{tikzpicture}[scale=0.9, every node/.style={font=\scriptsize}]
        \draw[accentdark, thick] (0,0) rectangle (5,2.8);
        \draw[attackcol, thick, dashed] (3,0) -- (3,2.8);
        \node at (1.5,2.45) {\textbf{0}=left};
        \node at (4,2.45) {\textbf{1}=right};
        \draw[accent, thick] (0,0) rectangle (3,2.8);
        \draw[attackcol, thick, dashed] (0,1.25) -- (3,1.25);
        \node at (0.8,1.9) {\textbf{0}=up};
        \node at (0.85,0.55) {\textbf{1}=down};
        \draw[goldcol, very thick] (0,0) rectangle (3,1.25);
        \fill[btccol] (1.45,0.62) circle (2pt);
        \node[btccol] at (2.25,0.62) {point};
      \end{tikzpicture}
    \end{column}
  \end{columns}
\end{frame}

% ── Predict-then-run ──────────────────────────────────────────────────────
\begin{frame}{Predict, then run}
  Before the demo, commit in \texttt{\#live-chat}:
  \begin{itemize}
    \item Given the bits \texttt{1 0}, which \emph{leaf} does the point land in?
    \item If we flip the first bit to \texttt{0}, does the point move
      \emph{left} or \emph{right}?
  \end{itemize}
  \vspace{6pt}
  \begin{center}\small
  Then we run it in the viewer and see who called it. Predicting first makes the
  mechanic \emph{stick}.
  \end{center}
\end{frame}

% ── Why a bijection ───────────────────────────────────────────────────────
\begin{frame}{Why it's genuinely a bijection}
  \begin{itemize}\small
    \item \textbf{Determinism.} Island discovery is driven by a PRNG seeded from
      the \emph{path so far} --- encode and decode walk the same path, find the
      same islands, pick the same splits.
    \item \textbf{No floating point.} Every bijection-critical step is
      integer / fixed-point (\textbf{I4F60}, 60 fractional bits) --- bit-identical
      on any machine. Floats' platform-dependent rounding would silently break the
      round-trip.
    \item \textbf{Clean partition.} Half-open intervals $[\min,\max)$ split a
      parent into two non-overlapping children --- every point in exactly one leaf.
  \end{itemize}
  \begin{pitfall}
    ``A fractal is chaotic, so this must be approximate.'' No --- the chaos is in
    the \emph{shape}; the \emph{arithmetic} is exact integer math. Same inputs,
    identical outputs, forever.
  \end{pitfall}
\end{frame}

\begin{frame}{Where does this lose people?}
  \gwfeedback{Three suspects: the \emph{splitting}, the \emph{determinism}
  (``why no floats''), or the \emph{``why a fractal at all.''} Tell us which one
  fogged over --- react A/B/C in \texttt{\#live-chat}.}
  \vspace{6pt}
  \begin{deeper}[Precision budget]
    Each bisection level costs $\sim$1 bit of positional precision. With 60
    fractional bits the encoder reliably places \textbf{32 bits per point} ---
    which is why bigger seeds use \emph{more stages}, not more bits per point.
  \end{deeper}
\end{frame}

% ── Hands-on ──────────────────────────────────────────────────────────────
\begin{frame}{Hands-on: round-trip a mini-seed}
  \begin{tryit}[In session]
    \begin{enumerate}\small
      \item Encode a \textbf{3-word (32-bit) mini-seed} in the viewer.
      \item Turn on area visualization (\texttt{V}) and step the bisection
        (\texttt{<} / \texttt{>}).
      \item \textbf{Decode the point back} and confirm you get the identical
        three words.
    \end{enumerate}
  \end{tryit}
  \vspace{2pt}
  \begin{center}\small\color{graytext}
  The skeptic's reassurance: your \emph{exact} words come back. Don't take our
  word --- watch them return.
  \end{center}
\end{frame}

% ── Check-out ─────────────────────────────────────────────────────────────
\begin{frame}{Check-out}
  \begin{columns}[T]
    \begin{column}{0.62\textwidth}
      \begin{block}{In \texttt{\#feedback}}
        \begin{enumerate}\small
          \item Muddiest point (name the moment / slide).
          \item Per-session form --- did the \emph{round-trip} convince you it's
            truly lossless?
          \item Did the hands-on \textbf{run on your machine}? If not, where?
        \end{enumerate}
      \end{block}
    \end{column}
    \begin{column}{0.36\textwidth}
      \memeslot{fractal-zoom.png}{a deep viewer zoom --- ``there's a whole world
      down here''}
    \end{column}
  \end{columns}
\end{frame}

\punch{We can encode. Next:\\[6pt] why nobody --- not even you ---\\[6pt] can rush it.}

\end{document}
